\section{Teknik Analisis Data}

Data sembilan metrik pada Tabel \ref{tab:definisi-operasional-metrik} yang diperoleh dari 30 replikasi setiap kombinasi algoritma–dataset–skenario dianalisis secara statistik untuk menguji H1–H4. 
Uji Shapiro–Wilk digunakan sebagai pemeriksaan awal distribusi data. 
Namun, karena kinerja metaheuristik bersifat stokastik dan tidak selalu berdistribusi normal, analisis utama menggunakan prosedur non-parametrik, sesuai praktik evaluasi metaheuristik \parencite{omran_permutation_2022}.

Untuk Eksperimen 1 dan 2 yang membandingkan tujuh algoritma, digunakan uji Friedman sebagai uji omnibus. 
Jika diperoleh perbedaan signifikan pada $\alpha=0{,}05$, analisis dilanjutkan dengan uji Wilcoxon \textit{signed-rank} antara CoDA-EDA dan setiap \textit{baseline}, disertai koreksi Holm–Bonferroni. 
Perbandingan dilakukan pada blok yang sama berdasarkan konfigurasi dataset, tingkat dependensi, dan replikasi yang berpasangan. 
Koreksi diterapkan per keluarga pengujian dalam satu dataset, yaitu enam perbandingan CoDA-EDA terhadap enam \textit{baseline}, dan tidak digabungkan lintas dataset.

Jumlah replikasi juga menentukan resolusi uji statistik. 
Pada uji Wilcoxon \textit{signed-rank} dua sisi tanpa \textit{tie} atau selisih nol, nilai-$p$\label{first:symbol-p} minimum untuk $n$ pasangan adalah $2/2^n$. Dengan lima pasangan, nilai minimum tersebut adalah $0,0625$ sehingga tidak dapat mencapai $\alpha=0{,}05$, sedangkan enam pasangan secara prinsip sudah memungkinkan. 
Oleh karena itu, 30 replikasi digunakan untuk menyediakan jumlah pengamatan yang memadai sekaligus menangkap variabilitas stokastik. 
\textit{Seed} pembangkitan dataset dan \textit{seed} algoritma dipisahkan agar variasi data tidak tercampur dengan variasi proses optimasi.

Eksperimen 3 dan 4 masing-masing membandingkan dua kondisi, yaitu CoDA-EDA terhadap iBOA serta \textit{warm-start} terhadap \textit{re-mining} penuh. 
Karena hanya terdapat dua kondisi, keduanya dianalisis langsung menggunakan uji Wilcoxon \textit{signed-rank} berpasangan tanpa uji omnibus.

Taraf signifikansi ditetapkan sebesar $\alpha=0{,}05$. Selain nilai-$p$, hasil dilengkapi ukuran efek non-parametrik, yaitu \textit{rank-biserial correlation} untuk uji Wilcoxon \parencite{peres_effect_2026}. 
Seluruh analisis dilakukan menggunakan pustaka SciPy pada Python.

Kualitas kebijakan dihitung terhadap \textit{log} yang digunakan dalam proses \textit{mining} sehingga F1, WSC, dan \textit{coverage} diinterpretasikan sebagai fidelitas rekonstruksi kebijakan, bukan kemampuan generalisasi prediktif terhadap permintaan yang belum diamati. 
Interpretasi ini sesuai dengan tujuan \textit{ABAC policy mining}, yaitu memperoleh kebijakan yang merepresentasikan pola keputusan pada \textit{log} akses.

Untuk hipotesis yang menyatakan kualitas sebanding atau non-inferior, pengujian tidak hanya didasarkan pada ketidakberhasilan menolak hipotesis nol. 
Margin non-inferioritas ditetapkan sebelum analisis, yaitu $\delta=0{,}02$ untuk F1 serta margin praktis yang ditetapkan untuk WSC. 
Non-inferioritas dievaluasi menggunakan selang kepercayaan 95\% terhadap selisih antarkondisi. 
Untuk F1, CoDA-EDA dinyatakan non-inferior apabila batas bawah selisih $F1_{\text{CoDA}}$-$F1_{\text{baseline}}$ berada di atas $-\delta$. 
Untuk WSC, karena nilai yang lebih rendah lebih baik, arah kriteria disesuaikan dengan definisi selisih yang digunakan.

Rancangan analisis statistik H1–H4 dirangkum pada Tabel \ref{tab:uji-statistik} berdasarkan eksperimen sumber data, kondisi yang dibandingkan, prosedur pengujian, dan ukuran efek yang digunakan.

\begin{longtable}{|>{\raggedright\arraybackslash}p{1.7cm}|>{\raggedright\arraybackslash}p{2.1cm}|>{\raggedright\arraybackslash}p{4cm}|>{\raggedright\arraybackslash}p{4cm}|>{\raggedright\arraybackslash}p{1.8cm}|}
\caption{Rancangan Uji Statistik per Hipotesis}
\label{tab:uji-statistik} \\
\hline
\textbf{Hipotesis} & \textbf{Eksperimen} & \textbf{Kondisi yang dibandingkan} & \textbf{Prosedur uji} & \textbf{Ukuran efek} \\
\hline
\endfirsthead

\multicolumn{5}{c}{{\bfseries \tablename\ \thetable{} -- lanjutan dari halaman sebelumnya}} \\
\hline
\textbf{Hipotesis} & \textbf{Eksperimen} & \textbf{Kondisi yang dibandingkan} & \textbf{Prosedur uji} & \textbf{Ukuran efek} \\
\hline
\endhead

\hline
\multicolumn{5}{r}{Lanjut ke halaman berikutnya} \\
\endfoot

\hline
\endlastfoot

% ==================== DATA BARIS ====================
H1 --- Kualitas kebijakan & Eksperimen 1 & Tujuh kondisi algoritma sekaligus, per dataset $\times$ tingkat korelasi & Friedman (omnibus) $\rightarrow$ post-hoc Wilcoxon signed-rank CoDA-EDA vs tiap baseline, koreksi Holm-Bonferroni & Rank-biserial \\
\hline
H2 --- Jejak memori & Eksperimen 2 & Tujuh kondisi algoritma sekaligus, per dataset & Friedman (omnibus) $\rightarrow$ post-hoc Wilcoxon signed-rank berkoreksi Holm-Bonferroni & Rank-biserial \\
\hline
H3 --- Waktu pembaruan & Eksperimen 3 & Dua kondisi (CoDA-EDA warm vs iBOA; serta warm vs always) & Wilcoxon signed-rank berpasangan langsung (tanpa omnibus) & Rank-biserial \\
\hline
H4 --- Warm-start & Eksperimen 4 & Dua kondisi (warm-start vs re-mining penuh) & Wilcoxon signed-rank berpasangan langsung (tanpa omnibus) & Rank-biserial \\
\hline

\end{longtable}
Tabel \ref{tab:uji-statistik} menjadi acuan pelaksanaan uji statistik atas hasil \textit{Demonstration}.
Dengan tersusunnya rancangan analisis ini, seluruh instrumen metodologi telah lengkap.